\section{Introduction}

Repeated observations impose different resolutions on different
perturbations of a quantum measurement. An interior change of an
eigenvalue is resolved on the scale $N^{-1/2}$, whereas a rotation of
a projective measurement can be resolved on the scale $N^{-1}$.
The two scales meet at singular effects and repeated eigenvalues.
We determine the resulting geometry on a compact matrix body, calculate
its description entropy, and construct a common learning procedure
whose loss is the distinguishability of all future experiments with
at most $N$ uses.

Fix a finite input dimension $d$ and consider
\[
 \mathfrak E_d=\{E=E^*\in M_d(\C):0\preceq E\preceq I_d\}.
\]
The effect $E$ specifies the memoryless ordered binary measurement
\[
 \mathcal M_E(\rho)=\operatorname{tr}(E\rho)|1\rangle\langle1|
       +\operatorname{tr}((I_d-E)\rho)|0\rangle\langle0|.
\]
The input is consumed and the output is classical. The distance $d_N$
is the supremum of the unhalved trace norm between the final states
of a common experiment with at most $N$ calls. The experiment may
entangle its inputs, retain an arbitrary reference and quantum memory,
use feedback, and stop publicly. Thus $0\le d_N\le2$; classical total
variation is one half of the corresponding classical trace norm.
The nonadaptive restriction $d_N^{\rm na}$ retains entangled input
blocks and references and omits feedback between calls.

\subsection{An intrinsic comparison and its entropy}
For $E,F\in\mathfrak E_d$ put
\[
 M=\tfrac12(E+F),\qquad H=F-E,\qquad V=M(I_d-M),
\]
and define the midpoint modulus
\[
 Q_N(E,F)^2=N\langle H,
       (L_V+R_V+N^{-1}\operatorname{Id})^{-1}H\rangle_{\rm HS}.
\]
The inverse acts on the real Hilbert space of Hermitian matrices.
Theorem~\ref{thm:matrixmetric76} proves, for every $d,N\ge1$,
\begin{equation}\label{eq:intrometric77}
 \frac{\min\{1,Q_N(E,F)\}}{8192d}
 \le d_N^{\rm na}(\mathcal M_E,\mathcal M_F)
 \le d_N(\mathcal M_E,\mathcal M_F)
 \le\min\{2,8Q_N(E,F)\}.
\end{equation}
No spectral margin or multiplicity assumption is imposed. In particular
$d_N\le65536d\,d_N^{\rm na}$. This is a fixed-dimensional comparison
of optimal values. The modulus $Q_N$ is neither identified with an
optimal value nor assumed to satisfy a triangle inequality.

A horizontal measurement dilation yields the adaptive tangent bound.
A regularized Sylvester equation separates an arbitrary tangent from a
remainder controlled by the hybrid inequality. Operator concavity of
$G(I_d-G)$ then integrates the tangent bound along a finite segment,
including boundary endpoints. The converse uses Bernoulli inputs and
exact two-dimensional compressions. All required qubit statements and
their proofs are included below.

Let $\mathcal C_N(\mathfrak E_d,\delta)$ denote the least number of
legal memoryless binary centres needed to cover the body in $d_N$.
For every fixed $d$, Theorem~\ref{thm:matrixcover76} gives
\begin{equation}\label{eq:introentropy77}
 \mathcal C_N(\mathfrak E_d,\delta)
 \asymp_d N^{d^2/2}[\log(N+2)]^{\lfloor d/2\rfloor}\delta^{-d^2},
 \qquad N\ge1,\quad0<\delta\le\delta_d.
\end{equation}
The positive cap $\delta_d$ and the comparison constants may depend
on $d$. The proof first bounds the weighted volume of actual
operational balls uniformly over every boundary stratum. The total
volume is then a spectral integral whose ordered endpoint sectors have
cumulative exponents $S_j^2/2$, where $S_j$ is a signed prefix sum.
Balanced prefixes and matching nested endpoint boxes produce the
logarithmic multiplicity $\lfloor d/2\rfloor$.

\subsection{Learning with a future operational loss}
The statistical problem supplies calls to an unknown effect, without
its matrix entries. For every fixed $d$, Theorem~\ref{thm:matrixlearning77}
constructs one learner, common to all $E\in\mathfrak E_d$, satisfying
\begin{equation}\label{eq:introlearning77}
 \sup_E\mathbb P_E\{d_N(\mathcal M_E,\mathcal M_{\widehat E})
                         >\delta\}\le\eta,
 \qquad
 M\le C_d N\delta^{-2}\log(1/\eta).
\end{equation}
Here $N\ge1$, $0<\delta\le\delta_0$ for an absolute $\delta_0>0$,
and $0<\eta\le1/8$.
The bound on the training calls $M$ holds on every record. Each
entangled block uses at most $N$ calls. A scalar Bernoulli subfamily
gives the matching lower order even for arbitrary reference-assisted
adaptive training. More precisely, the upper bound is
$Cd^4N\delta^{-2}\log(d/\eta)$ with an absolute $C$.
The dependence on $d$ is not asserted to be optimal.

The construction begins with a qubit learner for unknown bias,
contrast, and direction. Independent fair signs remove the bias from
GHZ parity observations. A dyadic amplitude search stops at its first
rejected scale, preserving a small-cone invariant; fresh observations
complete the direction and endpoint estimates. Its confidence ledger
controls both exceptional records and the maximum training cost.

For a general matrix effect, an additional calibration stage is
essential. It estimates the two support boundaries in relative
Loewner order at scale $N^{-1}$ and controls the absolute matrix error
at scale $N^{-1/2}$. In the calibrated eigenbasis, these two guarantees
control the extra variance introduced by every two-dimensional
compression. Qubit learning on those compressions then gives
compatible matrix coordinates in the regularized quadratic form.
Projection onto the legal effect body completes the estimate. The
calibration and compression budgets are summed conditionally on the
observed record, with no additional $\log N$ factor.

\subsection{Independent training blocks and dimension}
The maximal size of an independent quantum experiment is a second
training resource.  At the boundary of a block, a learner may jointly
process all completed outputs and choose its next preparation from the
classical record.  The next block, including its reference and unused
inputs, is isolated from earlier quantum registers until its last call.
This permits collective final estimation and classical feedback.
It excludes coherent transfer from older quantum memory into a block
whose queries are not yet complete.

Let $1\le b\le N$ be the maximal declared block length.
For every fixed $d\ge2$, Theorem~\ref{thm:blocklearning78} proves
\begin{equation}\label{eq:introblock78}
 M_b^\star(d,N,\delta,\eta)
 =\Theta_d\!\left(\frac{N^2}{b\delta^2}\log\frac1\eta\right),
\end{equation}
for an absolute small-error range and $0<\eta\le1/8$.
The explicit upper bound is
\[
 Cd^4\frac{N^2}{b}\delta^{-2}\log(d/\eta).
\]
A horizon-blocking inequality gives the common learner, while a
history-weighted root-fidelity bound gives the converse even with
arbitrary quantum storage and collective processing at block boundaries.
In dimension one the device is a coin and the cost remains
$\Theta(N\delta^{-2}\log(1/\eta))$, independently of $b$.

There is also a joint dimension obstruction.
Theorem~\ref{thm:dimensionlower78} proves
$M\ge c d^2N\delta^{-2}$, already on
$I_d/4\preceq E\preceq3I_d/4$, under arbitrary coherent adaptive
training.  If every candidate effect lies within operator radius $r$
of $I/2$, one query increases information about its finite unknown
label by at most $4r^2$.  A matrix packing and Fano's inequality
then give the lower bound.
In this fixed interior the operational distance has a dimension-free
comparison with $\min\{1,\sqrt N\|E-F\|_{\rm op}\}$.
Reanalysing the binary Choi estimator of \cite{MB67} consequently gives
\begin{equation}\label{eq:introinterior78}
 M_{\rm int}^\star(d,N,\delta,1/8)
 =\Theta(d^2N\delta^{-2})
\end{equation}
with absolute constants, achieved using independent one-call probes.
At $N=1$ this also gives the joint dimension--accuracy order
$\Theta(d^2\epsilon^{-2})$ for binary operator-norm learning.
Thus existing independent-probe acquisition does admit a stronger
future-loss analysis in the interior, whereas the projective boundary
in \eqref{eq:introblock78} prevents that rate uniformly on the full body.

\subsection{Confidence as a joint resource}
The fixed-confidence law has a refinement uniform in all four
parameters.  For $d,N\ge1$, $0<\delta\le2^{-13}$ and
$0<\eta\le1/8$, Theorem~\ref{thm:interiorconfidence80} proves
\begin{equation}\label{eq:introconfidence80}
 M_{\rm int}^{\star}(d,N,\delta,\eta)
 =\Theta\!\left(\frac{N}{\delta^2}
                   \left(d^2+d\log\frac1\eta\right)\right).
\end{equation}
The constants are absolute.  The upper procedure uses independent
one-call Choi acquisitions, and the lower permits arbitrary coherent
adaptive training.  Thus the same order holds under every fresh-block
cap $1\le b\le N$.  At $N=1$, it also gives the sharp joint
dimension, accuracy and confidence law for binary operator-norm
learning.

For the confidence converse, compare the central coin $I_d/2$ with
$d$ diagonal single-coordinate perturbations.  A revealed basis
outcome permits a relative-entropy chain rule through arbitrary
retained quantum memory.  Under the central hypothesis, the sum of
the expected occupation counts is the entire training budget.
Binary testing consequently forces a $d\log(1/\eta)$ contribution.
Combining this with the matrix-packing lower bound gives the sum in
\eqref{eq:introconfidence80}; no product of separate lower bounds
is taken.  For the upper bound, the confidence-dependent estimator
of \cite[Corollary~III.9]{MB67} is run at a rational failure
probability exponentially small in $d$.  Independent repetition
and selection of a metric majority extend its guarantee to every
$\eta$ at the same additive cost.

The same estimator gives a further upper bound on the complete effect
body.  At operator accuracy $\delta/(2N)$, the hybrid inequality gives
future error at most $\delta$.  For any fresh-block cap, one may
therefore select between this one-call construction and the block
learner, obtaining
\[
 M_b^\star\le C\frac{N^2}{\delta^2}
 \min\left\{d^2+d\log(1/\eta),\frac{d^4}{b}\log(d/\eta)\right\}
\]
in their common range; see \eqref{eq:fullbodyconfidenceupper80}.
Together with \eqref{eq:fullbodyconfidencelower80}, this displays
the current joint full-body bounds.

\subsection{Deterministic descriptions and resources}
Theorem~\ref{thm:matrixcodec77} gives a deterministic finite rational
codebook at the order in \eqref{eq:introentropy77}. Its greedy rule
uses exact comparisons of $Q_N^2$ on a finite legal matrix grid.
The operational separation in \eqref{eq:intrometric77} controls its
cardinality, and an inward rounding map gives coverage at every rank.
The encoder and decoder reproduce the same dictionary from the public
parameters, with no supplied custom dictionary. Every fixed-length
word decodes to a legal effect. The required payload is
\[
 \frac{d^2}{2}\log_2N+
 \lfloor d/2\rfloor\log_2\log(N+2)
             +d^2\log_2(1/\delta)+O_d(1).
\]
The construction is finite and exact; its enumeration cost is large.
Optimal payload length does not imply efficient classical computation.
The direct qubit code remains available with its own detailed resource
accounting. Corollary~\ref{cor:matrixlearnedcode77} combines the common
learner with this exact encoder, attaining the learning and payload
orders simultaneously without additional calls to the unknown device.
Corollary~\ref{cor:blocklearnedcode78} preserves this conclusion under
a prescribed independent-block cap.  The finite-control theorem
approximates all trusted preparations of our common learners with a
pathwise error budget.  The original loss, confidence, and call orders
are retained; classical descriptions and trusted realization errors
are charged explicitly.


\subsection{Dimension-uniform operational descriptions}
The interior learning result has a matching description theory.
For $\mathfrak I_d=[I_d/4,3I_d/4]$, all $d,N\ge1$ and
$0<\delta\le2^{-13}$, Theorem~\ref{thm:interiorentropy79} gives
\begin{equation}\label{eq:introinteriorentropy79}
 \left(\frac{\sqrt N}{2048\delta}\right)^{d^2}
 \le\mathcal C_N(\mathfrak I_d,\delta)
 \le\left(\frac{25\sqrt N}{\delta}\right)^{d^2}.
\end{equation}
The constants are absolute and the lower bound allows arbitrary
legal binary centres.  Theorem~\ref{thm:interiorcodec79} realizes
the upper bound by finite exact rational comparison and public
reconstruction.  Unlike entrywise transmission, its index does not
incur a $d^2\log d$ penalty.  Its length is
\[
 A_d(N,\delta)+O(d^2),\qquad
 A_d(N,\delta)=\tfrac12d^2\log_2N+d^2\log_2(1/\delta),
\]
uniformly in dimension, horizon and accuracy.

If shared randomness and a worst-target failure probability $\eta$
are allowed, Theorem~\ref{thm:prefixrate79} determines the optimal
expected one-way prefix length as
$(1-\eta)A_d(N,\delta)+O(d^2)$, uniformly for $0\le\eta\le1/8$.
The lower bound allows a randomized decoder and counts all nonpublic
information.  The upper bound uses a public target-independent
abort event; its flag is not a hidden transmission.
Finally, Theorem~\ref{thm:interiorlearnedcode79} gives one common
finite-output learner with the joint optimal orders
$M=\Theta(d^2N\delta^{-2})$ and $B=A_d(N,\delta)+O(d^2)$ at
failure $1/8$, using independent one-call acquisitions.  This is a
query and communication result; trusted tomography synthesis and
exhaustive decoder computation are separate costs.
The absence of the endpoint logarithm in
\eqref{eq:introinteriorentropy79} identifies the different interior
law without changing the retained complete-body theorem.

\subsection{An effective finite collective learner}
The ideal finite-output learning theorem has an effective counterpart.
For rational $\delta,\eta$ in the preceding range,
Theorem~\ref{thm:effectiveinterior80} computes an entire finite
classical specification before the first call to the unknown device.
It retains the query order in \eqref{eq:introconfidence80} and
the index length $A_d(N,\delta)+O(d^2)$.

The construction fixes the public rational dictionary and asks for
a finite collective readout whose uniform operator-error risk is
sufficient for the operational guarantee.  The unnormalized Choi
record is polynomial in the unknown effect.  Readout positivity,
normalization and its uniform success condition therefore form a
first-order formula over the real numbers with rational coefficients.
Effective quantifier elimination decides feasibility, and an
algebraic-point selection gives exact finite readout data.
The existence theorem bounds the first feasible call count; its
continuous measurement operators are not supplied to the search.
Only the resulting finite readout is realized.  The pathwise
trusted-control estimate preserves the loss and confidence with
finite precision.  Quantifier elimination, reference storage,
known-control synthesis and public dictionary reconstruction remain
charged resources; the theorem makes no efficient computation claim.

\subsection{Prior work and organization}
Horizontal Kraus gauges, channel-extension metrology, and integration
of local distinguishability have established antecedents
\cite{FI76,DKG67,DM76,KGAD76,YF76,SD76}.
Inverse-Sylvester geometry and spectral matrix integration are likewise
classical \cite{Dittmann76,SZGeom76}. The explicit closed-body
midpoint comparison, its nested endpoint entropy, and learning under
the future operational loss are the theorem objects studied here.

Mele--Bittel \cite{MB67} and
Zambrano--Ramos-Calderer--Kueng \cite{ZRK76} already provide common
estimators for unknown measurements in one-use losses. Their access,
rank, dimension, and confidence statements are compared at theorem
level in Section~\ref{sec:matrixcomparison76}. The distinction in
\eqref{eq:introlearning77} is the dependence on the future horizon
under the anisotropic loss $d_N$. The phase-estimation components
have antecedents in \cite{HB76,KLY76}; the conditional construction
used here is proved directly.

The matrix comparison and entropy are followed by the qubit learning
construction, its matrix extension, and the exact codebook theorem.
The block-resource law, joint dimension bounds, uniform interior descriptions,
confidence refinement, and effective finite-control construction then identify
how acquisition, confidence and preparation affect learning. The later qubit geometry and the final prerequisite
appendix make the focused article self-contained. The global covering law and the matching description order use fixed
dimension and small error.  The full-body bounds compare two explicit
constructions and their joint lower obstructions, while the sharp joint law is
proved on the stated interior jointly in dimension, horizon, accuracy and confidence.  The exact matrix
code is a finite exhaustive construction.  All results concern ordered
binary, input-consuming, classical-output measurements.
